\par\noindent\textbf{(1)} 单个位点占有数 \(n=0,1\)，权重为 \(1\) 与 \(\exp[\beta(\mu+\epsilon)]\)，故 \(\xi_1=1+e^{\beta(\mu+\epsilon)}\)，整表面 \(\Xi=\xi_1^M\)。这里位点独立且每位点不可重复占据，不能写成任意玻色占有的几何级数。
\par\noindent\textbf{(2)} 平均覆盖率 \(\theta=\langle n\rangle=(\beta M)^{-1}\partial_\mu\ln\Xi=1/[1+e^{-\beta(\mu+\epsilon)}]\)。对于稀薄理想气体库 \(e^{\beta\mu}\propto p\)，可写 \(\theta=K(T)p/[1+K(T)p]\)，其中平衡常数含吸附束缚因子 \(e^{\beta\epsilon}\)。
\par\noindent\textbf{(3)} 每位点为Bernoulli变量，故 \(\operatorname{Var}(N_a)=M\theta(1-\theta)\)。直接微分 \(\partial_\mu\langle N_a\rangle=\beta M\theta(1-\theta)\)，验证巨正则涨落恒等式 \(\operatorname{Var}(N_a)=k_BT\,\partial_\mu\langle N_a\rangle\)。
\par\noindent\textbf{(4)} 在 \(Kp\ll1\) 时 \(\theta\simeq Kp\)，与压力线性；在 \(Kp\gg1\) 时 \(\theta\to1\)，最大吸附数为 \(M\)。若吸附分子之间存在侧向相互作用，位点独立假设破坏并可能产生协同吸附，不能用Langmuir曲线拟合临界覆盖行为。
